\section{Rich lines and overlapping hyperplanes}\label{sec:geometry}

The score used by the predictor is a sum over hyperplanes through a query
point.  Its terms can depend strongly on one another when two hyperplanes
contain many common training points.  We now bound these overlaps.  The
training-set construction in Lemma~\ref{lem:training} removes large flat
concentrations from each query's samples; the remaining concentrations are
controlled by a polynomial argument for lines and a counting argument for
planes.

Low-degree polynomials vanishing on point sets play a central role in
Dvir's finite-field Kakeya argument \cite{Dvir2009}.
The rich-line proof here uses the polynomial interpolation and tangent-plane
methods of Guth and Katz~\cite{GuthKatz2010} and Elekes, Kaplan, and
Sharir~\cite{ElekesKaplanSharir2011}.  The positive-characteristic
obstruction to these methods is treated by Ellenberg and
Hablicsek~\cite{EllenbergHablicsek2016}.  Here the defining polynomial
has degree below the prime characteristic, and we give the required
algebraic argument in full.
Koll\'ar's positive-characteristic incidence estimate
\cite[Theorem~2]{Kollar} also motivated the rich-line framework. The
plane-capped estimate below is proved directly and does not invoke that
theorem as a premise.

Throughout this section, $q$ is a sufficiently large prime,
$\sigma=\log q$, $\log\sigma=o(P)$ and $P=o(\sigma)$.  Write
\[
 \begin{gathered}
 n=q^{1+v},\qquad \frac n4\le n'\le n,\qquad
 g_0=(v-\tfrac12)\sigma,\qquad g_1=(v-1)\sigma,\\
 \chi=\frac P{100},\qquad \chi_0=\frac P{10000}.
 \end{gathered}
\]
All constants below are absolute.  Bounds denoted by $o(n)$ are uniform
under the parameter assumptions just stated and those of the relevant
lemma.

\subsection{A polynomial bound for rich lines}

We first prove the only algebraic ingredient.  Its characteristic
restriction will enter through a polynomial whose degree is less than $q$.

\begin{lemma}[Rich lines with a plane cap]\label{lem:rich-lines}
Suppose $n=q^{1+v}$, $v\le 3/2+O(1/\sigma)$,
$n/4\le n'\le n$, $g_0>\chi_0$, and
$e^{-g_0/20}<a\le2$.  Put
\[
 M=\frac{n'a}{q},\qquad
 K=n^{4/3}q^{-1}e^{-g_0/5}.
\]
Let $X$ be a set of at most $n$ points of $\mathrm{PG}(d,q)$, where
$3\le d\le4$, such that every projective plane contains at most $K$
points of $X$.  The number of lines containing at least $M$ points of $X$
is at most $Cn/M$.
\end{lemma}

\begin{proof}
The inequalities that make the argument work are
\begin{equation}\label{eq:rich-line-parameters}
 \frac{M}{n^{1/3}}\ge c e^{37g_0/60},\qquad
 \frac{M^2}{K}\ge c e^{23g_0/30},\qquad
 cq\le\frac nM\le Cq e^{g_0/20}.
\end{equation}
They follow by substituting $v=1/2+g_0/\sigma$ and using
$a>e^{-g_0/20}$ and $n'\ge n/4$.  In particular, their first two
ratios tend to infinity.

Work temporarily over an algebraic closure of $\mathbb F_q$.  A generic
linear map to a four-dimensional vector space preserves the independence
of every independent subset of at most four points of $X$.  Indeed, each
of these finitely many requirements defines a nonempty open condition on
the map, and the field is infinite.  The resulting projection to
projective three-space preserves distinct points and distinct rich lines.
It also preserves the plane cap: four projected points in a plane cannot
have come from four independent points.  Choose an affine chart containing
all projected points.

Sample these points independently with probability
\[
 f=A\frac{\sqrt n}{M^{3/2}},
\]
where $A$ is a sufficiently large fixed constant.  By
\eqref{eq:rich-line-parameters}, $f=o(1)$ and
$fM=A\sqrt{n/M}\ge cA\sqrt q$.  There are at most $n^2$ lines containing
at least two points of $X$.  Chernoff's inequality and a union bound imply
that, with probability $1-o(1)$, every rich line contains more than $fM/2$
sampled points.  Markov's inequality gives probability at least $1/2$ that
the total sample size is at most $2fn$.  Choose a sample satisfying both
conditions.

Counting monomials in three variables produces a nonzero polynomial
vanishing on that sample, of degree
\[
 D_0\le C A^{1/3}\sqrt{n/M}.
\]
For sufficiently large fixed $A$, this degree is less than $fM/2$, so the
polynomial vanishes identically on every rich line.  Replacing it by its
squarefree part does not increase the degree or lose any such line.  Also
\begin{equation}\label{eq:degree-characteristic}
 D_0=o(M),\qquad D_0=O(q^{21/40})<q,
\end{equation}
because $g_0\le\sigma+O(1)$.

First consider its plane factors.  At most $D_0^2$ lines lie in two
different factor planes.  Fix one factor plane containing $r$ of the
remaining lines, and choose $M'=\lfloor M\rfloor$ points on each line.
If $d_z$ is their multiplicity at a selected point, the plane cap and the
fact that two lines meet at most once give
\[
 (rM')^2\le K\sum_z d_z^2
 \le K\bigl(rM'+r(r-1)\bigr).
\]
Consequently $r\le KM'/((M')^2-K)=o(M')$.  Delete from each line its
intersections with all other factor planes, removing at most $D_0$ points.
The remaining union in this plane has size at least
\[
 r(M'-D_0)-\binom r2\ge \frac{rM'}2.
\]
These unions are disjoint for different factor planes.  They account for
at most $Cn/M$ lines in total.

Let $F_0$ be the squarefree product of the nonplane factors, and consider
the remaining rich lines.  Call a point of $X$ busy if it lies on at least
three of these lines.  At a busy point where $\nabla F_0\ne0$, the three
distinct line directions lie in the tangent plane, and the Hessian
quadratic form vanishes on them.  A quadratic form on a two-dimensional
space with three distinct projective zeros is identically zero.  Since
the characteristic is odd, polarization shows that the Hessian bilinear
form vanishes on that tangent plane.  Thus every busy point is a zero of
all the polynomials
\begin{equation}\label{eq:hessian-polynomials}
 Q_{ij}=(e_i\mathbin\times\nabla F_0)^t
              \operatorname{Hess}(F_0)
              (e_j\mathbin\times\nabla F_0),\qquad 1\le i,j\le3.
\end{equation}
This is also true at a point with zero gradient.  There are at most $2n$
nonbusy point--line incidences.  Except for $Cn/M$ lines, every remaining
line therefore has at least $M'/2$ busy points.  As $\deg Q_{ij}<3D_0$ and
$M'/2>3D_0$ for large $q$, each such line lies in every $Q_{ij}$.

We claim that no irreducible factor $F$ of $F_0$ divides all the
$Q_{ij}$.  Since $\deg F<q$, some derivative of $F$, say $F_z$, is
nonzero.  Squarefreeness implies that $F$ and $(F_0)_z$ are coprime in
$\overline{\mathbb F}_q(x,y)[z]$.  Clearing a B\'ezout identity in that
field and specializing $x,y$ away from its denominators and the leading
coefficient of $F$ produces a point on $F$ where $(F_0)_z\ne0$.
After translation, the formal implicit-function theorem writes the
surface locally as $z=h(x,y)$.

If every $Q_{ij}$ vanished on this surface, the Hessian bilinear form of
$F_0$ would vanish on its tangent space: the vectors
$e_i\mathbin\times\nabla F_0$ span that space.  Differentiating
$F_0(x,y,h(x,y))=0$ twice would then give
\[
 h_{xx}=h_{xy}=h_{yy}=0.
\]
For a nonzero monomial $x^u y^w$ of $h$, these identities force
$(u,w)$ modulo the prime $q$ to be one of
$(0,0),(1,0),(0,1)$.  Thus $h$ agrees with its affine-linear part $h_1$
through total degree $q-1$.  The polynomial $F_0(x,y,h_1(x,y))$ has degree
less than $q$ and vanishes to order at least $q$, so it is zero.  This
would give a plane factor of $F_0$, a contradiction.

Over the algebraic closure, a constant linear combination $Q$ of the
$Q_{ij}$ can therefore be chosen coprime to $F_0$: avoid the finitely many
proper linear subspaces of coefficients corresponding to its irreducible
factors.  We use the following elementary consequence of resultants.
\begin{quote}
Two coprime polynomials in affine three-space, of degrees $D_1,D_2$,
contain at most $2D_1D_2$ common lines.
\end{quote}
Here is a proof sufficient for this application.  Choose a linear
coordinate $z$ nonconstant on every line in any fixed finite collection
of common lines.  For a generic value $s$, the restrictions to $z=s$
remain nonzero and coprime.  To see this, use the nonzero resultants
eliminating $y$ and eliminating $x$, viewed over the respective fraction
fields; preserve a nonzero coefficient of each applicable resultant
under specialization.  A common factor of positive degree in either
variable would make the corresponding resultant vanish.  If both
polynomials are independent of an elimination variable, omit that
resultant.  Nonvanishing of the specialized polynomials is a further
open condition.  This argument remains valid when degrees drop: a common
factor supplies a kernel vector for the Sylvester map within the original
degree bounds.

Avoid also the finitely many values $s$ at intersections of selected
lines.  Their slice points are distinct.  Choose coordinates in the slice
so their first coordinates are distinct.  The resultant eliminating the
second coordinate is nonzero, has degree at most $2D_1D_2$, and vanishes
at each of those first coordinates.  If both polynomials are independent
of the second coordinate, coprimality gives no common point instead.
This proves the assertion.  Applied to $F_0,Q$, it bounds the remaining
lines by $O(D_0^2)=O(n/M)$, completing the proof.
\end{proof}

\subsection{The overlap estimate}

We now return to the training set.  To keep the geometric statement
independent of the score, the hyperplanes below can be any family of
bounded training density.

\begin{lemma}[Geometry of the training set]\label{lem:overlap-geometry}
Let $3\le d\le4$, $100qP<n=q^{1+v}\le Cq^{(d+1)/2}$, and let
$S'\subset\mathrm{PG}(d,q)$ be the training set in the small-cell case of
Lemma~\ref{lem:training}, with $n'=|S'|\in[n/4,n]$.
For each query point $x$, let $O_x$ be the union of its earliest containing
assignment cells in the captured hyperplane and plane lists; an absent
list contributes the empty set.  Put
\[
 X_x=S'\setminus(O_x\cup\{x\}),\qquad B=\frac{q^d}{n'},\qquad
 a_\ell=\frac{q|X_x\cap\ell|}{n'}
\]
for lines $\ell$ through $x$.
For each $x$, let $\mathcal H_x$ be any family of hyperplanes through $x$
such that $q|X_x\cap H|/n'\le2$ for every $H\in\mathcal H_x$.
For distinct members set
\[
 \theta_{HH'}=\frac{q|X_x\cap H\cap H'|}{n'},
\]
and let $E_x(a)$ count ordered distinct pairs with $\theta_{HH'}\ge a$.
Let $\Delta_x(a)$ be the maximum degree of this relation.

There is a set of $o(n)$ query points outside which, simultaneously for
all dyadic $a\in[q^{-10},2]$,
\begin{align}
 \#\{\ell\ni x:a_\ell\ge a\}
   &\le q^{2-2v}e^\chi a^{-100},\label{eq:local-rich-lines}\\
 E_x(a)&\le B^2e^{2\chi}a^{-200},\label{eq:overlap-pairs}\\
 \Delta_x(a)&\le Be^{2\chi}a^{-200}.\label{eq:overlap-degree}
\end{align}
The exceptional set is independent of the choice of the families
$\mathcal H_x$.

Let $p$ be an even integer in
$[10000\sigma/P,10000\sigma/P+2]$ and put $t=1/(100p)$.
After excluding at most $n\exp(CP)$ additional points,
\begin{equation}\label{eq:ambient-overlap-degree}
 \Delta_x(t)\le B e^{-3P}.
\end{equation}
These additional points are needed only for the ambient bound, and are
not included in the $o(n)$ validation exception.

Two further estimates hold without point exclusions.  Every line strength
is at most $Ce^{-g_1}$.  If $g_0>\chi_0$ and
$e^{-g_0/20}<a\le2$, the number of pairs $(x,\ell)$ with $x\in\ell$ and
$a_\ell\ge a$ is at most $Cq^2/a^2$.  For all strengths, the number of
possible point--line pairs is at most $Cq^7$.
\end{lemma}

\begin{proof}
We first prove the line estimates.  Since the radial lines partition
$X_x$, their number of strength at least $a$ is at most $q/a$.  This
implies \eqref{eq:local-rich-lines} if $g_0\le\chi_0$ or
$a\le e^{-g_0/20}$, because the ratio of the former bound to the latter
is
\[
 e^{2g_0-\chi}a^{99}=o(1).
\]
Suppose otherwise, and put $M=n'a/q$.  The plane-peeling step was then
performed, with threshold
$K_p=n^{4/3}q^{-1}e^{-g_0/5}$ and list length
$J_p\le n/K_p$.  In addition to
\eqref{eq:rich-line-parameters}, we have
\begin{equation}\label{eq:plane-list-length}
 J_p\le q^{1/2}e^{-2g_0/15},\qquad J_p/M\le C e^{-13g_0/12}=o(1).
\end{equation}

If the capped residual was retained, Lemma~\ref{lem:rich-lines} gives
$Cn/M\le Cq/a$ rich lines.  Each has $q+1$ rational centers, so there
are at most $Cq^2/a^2$ relevant center--line pairs.

If a captured prefix of planes was retained, a line not contained in any
listed plane has at most $J_p=o(M)$ training points.  On a line first
contained in plane $j$, earlier planes contribute at most $J_p$ points,
and later assignment cells contribute none.  Cell $j$ must therefore
contribute at least $M/2$ points.  For these points to survive the removal
of $O_x$, the center $x$ must belong to an earlier listed plane.  None of
those earlier planes contains the line, so at most $j$ centers qualify.
Write $f_j$ for the fraction of $S'$ in cell $j$.

If $a>Cf_j$, the incidence variance in that projective plane gives at
most $Cq^{2-v}f_j/a^2$ rich lines in it.  Multiplying by $j$ and summing
uses $\sum_j f_j\le1$ and $J_pq^{2-v}\le q^2$, giving $Cq^2/a^2$
center--line pairs.  If $a\le Cf_j$, the nonincreasing cell sizes from
the greedy plane peel imply $j\le C/a$.  There are at most $Cq^2$ lines
in each such plane, and summing the center multiplicities $j$ again gives
$Cq^2/a^2$.  Additional removal of a hyperplane cell can only decrease
these counts.

We have proved the asserted global center--line bound.  Markov's
inequality at the threshold in \eqref{eq:local-rich-lines} leaves a
fraction at most
\begin{equation}\label{eq:line-exception-fraction}
 C e^{g_1-\chi}a^{98}
\end{equation}
of $n$ exceptional centers.  If $g_1\le\chi_0$, this tends to zero even
after summing over $O(\sigma)$ dyads.  Otherwise every line has strength
at most $q(q+1)/n'\le Ce^{-g_1}$, so the same sum is still smaller.
The crude $Cq^7$ bound counts all points and all lines through them in
projective four-space; it also covers dimension three.

This proves the radial-line bounds. We next estimate the overlap edge
count $E_x(a)$ and maximum degree $\Delta_x(a)$.
Set $M=n'a/q$ for every strength
threshold $a$, including those covered by the trivial line bound.
In dimension three the intersection of two distinct hyperplanes through
$x$ is a radial line.  Each such line gives at most $Cq^2$ ordered pairs.
A fixed $H\in\mathcal H_x$ contains at most $2/a$ radial lines of
strength at least $a$, and each lies in at most $Cq$ other hyperplanes.
Hence
\[
 E_x(a)\le Cq^2\#\{\ell\ni x:a_\ell\ge a\},\qquad
 \Delta_x(a)\le Cq/a.
\]
Because $B\asymp q^{2-v}$ and $v\le1+O(1/\sigma)$, these prove
\eqref{eq:overlap-pairs} and \eqref{eq:overlap-degree} in dimension three.

It remains to analyze dimension four.  Intersections are now planes
through $x$.  Call an overlap plane dominant if one radial line carries
at least half its outside-own points, and nondominant otherwise.  A
dominant overlap of strength at least $a$ contains a line of strength at
least $a/2$.  Each such line gives at most $Cq^4$ ordered hyperplane pairs.
For a fixed $H\in\mathcal H_x$, the corresponding degree is at
most $Cq^2/a$.  Dividing this degree by $B\asymp q^{3-v}$ gives
$Ce^{g_1}/a$.  If $g_1\le\chi_0$, this meets
\eqref{eq:overlap-degree}; if $g_1>\chi_0$, a dominant overlap requires
$a\le Ce^{-g_1}$, and the same ratio is at most $C/a^2$.
Equation~\eqref{eq:local-rich-lines} also gives the required pair count.

In a nondominant plane of strength at least $a$, there are at least
$cM^2$ pairs of training points on distinct radial lines.  Each pair
together with $x$ determines the plane.  Thus the number of these planes
through $x$ is at most $Cq^2/a^2$.  Within a fixed $H\in\mathcal H_x$,
it is at most $C/a^2$, since $|X_x\cap H|\le2n'/q$.
Every plane lies in at most $Cq$ hyperplanes, so
\begin{equation}\label{eq:nondominant-crude}
 E_x(a)\le Cq^4/a^2,\qquad \Delta_x(a)\le Cq/a^2.
\end{equation}
The degree bound is always sufficient.  The pair bound is sufficient if
$g_1\le\chi_0$ or $a\le e^{-g_1/10}$: after dividing by $B^2$, its
extra factor is $Ce^{2g_1}/a^2$, which is then absorbed by
$e^{2\chi}a^{-200}$.

\paragraph{Rich planes in the remaining range.}
Suppose $g_1>\chi_0$ and $a>e^{-g_1/10}$.  Here
\begin{equation}\label{eq:rich-plane-parameters}
 M\gg q,\qquad M^2\gg n.
\end{equation}
If a captured hyperplane list was retained, take as its head the prefix
whose original cell sizes, before plane peeling, exceed $C_0q^2/a$,
where $C_0$ is a sufficiently large absolute constant.  Its length is at
most $na/(C_0q^2)$.  The remaining training tail has hyperplane cap
$C_0q^2/a$, by the maximal choice at every greedy step.  If instead the
hyperplane-capped residual was retained, use an empty head and the whole
training set as tail.

A plane not contained in a head hyperplane receives fewer than $M/4$
head points: each intersection is a line of at most $q+1$ points, and
the displayed head-length bound applies.  Let $\mathcal F$ be the global
family of planes containing at least $M/2$ tail points, and let
\[
 r=|\mathcal F|,\qquad u=q^{4-2v}a^{-3}.
\]
A line is contained in at most $Cn/M=O(u)$ members of $\mathcal F$,
since they are disjoint off that line and $M\gg q$.  A hyperplane is
a projective three-space and has at most $C_0q^2/a$ tail
points.  The incidence variance there bounds the number of its planes in
$\mathcal F$ by
\[
 \frac{Cq^2(C_0q^2/a)}{M^2}=O(u).
\]
The use of variance is valid because the mean is $O(q/a)=o(M)$.

For a query $x$, write $d_x$ for the number of members of $\mathcal F$
through $x$.  In the quotient by $x$ these become lines of a projective
three-space, with at most $Cu$ in any point or plane.  If $d_x>C'u$, a
constant fraction of pairs of these quotient lines are skew.  Indeed,
if no line meets half the others, this is immediate.  Otherwise fix a
line meeting at least half.  Among its neighbors, a given line has only
$O(u)$ companions with the same intersection point or the same spanning
plane with the fixed line.  All remaining pairs of neighbors are skew.
Two corresponding original planes intersect exactly in $x$.  Therefore
\begin{equation}\label{eq:rich-plane-second-moment}
 \sum_{d_x>C'u}d_x^2\le Cr^2.
\end{equation}
Count incidences at the tail points, splitting according to this
threshold and applying Cauchy--Schwarz:
\[
 rM/2\le Cnu+C\sqrt n\,r.
\]
By \eqref{eq:rich-plane-parameters}, $r\le Cnu/M$.  Equation
\eqref{eq:rich-plane-second-moment} now implies that
\begin{equation}\label{eq:rich-plane-local}
 d_x\le q^{4-2v}e^\chi a^{-4}
\end{equation}
outside at most $Cq^2e^{-2\chi}$ centers.

Next consider a rich plane contained in a head hyperplane, and take the
earliest such hyperplane, with index $j$.  Earlier noncontaining cells
contribute fewer than $M/4$ points, and later cells contribute none.
Richness therefore requires at least $M/2$ points of cell $j$.
For these to remain outside $O_x$, the center must lie in an earlier
hyperplane.  This allows at most $Cjq$ centers on the plane.
Let $f_j$ be the original cell size divided by $n$.  If $a>Cf_j$,
variance within hyperplane $j$ bounds these planes by $CBf_j/a^2$.
Summing their center multiplicities gives at most $Cq^3/a$, because
\[
 \sum_j j f_j\le \frac{na}{C_0q^2}\sum_j f_j
 \le\frac{na}{C_0q^2}.
\]
Markov's inequality at the threshold in \eqref{eq:rich-plane-local}
makes the exceptional fraction at most $Cq^{v-2}e^{-\chi}a^3=o(1)$.

If $a\le Cf_j$, then $x$ belongs to two original hyperplanes whose cell
fractions are at least $ca$, by the nonincreasing order.  There are at
most $C/a$ such hyperplanes; the union of their pairwise intersections
has at most $Cq^2/a^2$ points.  Relative to $n$, this is at most
$Ce^{-.8g_1}=o(1)$ in the present range.  All exceptions remain $o(n)$
after summing over the dyads.

We have bounded the number of rich overlap planes through every
remaining $x$ by $Cq^{4-2v}e^\chi a^{-4}$.  Each gives at most $Cq^2$
ordered hyperplane pairs.  This proves
\eqref{eq:overlap-pairs} in the remaining case.  The degree estimate
continues to follow from \eqref{eq:nondominant-crude}.  The fixed powers
of two from neighboring dyads and all constant factors are absorbed by
the exponential slack.  This completes the bounds outside $o(n)$ points.

\paragraph{The stronger bound used only for ambient points.}
If $g_1<-4P$, the dominant-line degree divided by $B$ is at most
$Ce^{g_1}/t\le e^{-3P}$.  If $g_1>\chi_0$, which can occur only in
dimension four, the maximum line strength $Ce^{-g_1}$ is smaller than
$t/2$; there is no dominant strong overlap.

In the remaining range $-4P\le g_1\le\chi_0$, exclude every center with
an outside-own radial line of strength greater than $t/2$.  Here
$g_0=\sigma/2+g_1>\chi_0$ and
$t/2\gg e^{-g_0/20}$, so the global center--line bound already proved
costs at most
\[
 Cq^2/t^2\le n e^{CP}
\]
centers.  There are now no strong pairs in dimension three.  In dimension
four only nondominant strong pairs remain, and
\[
 \Delta_x(t)/B\le Cq^{v-2}t^{-2}
 \le Cq^{-1/2+o(1)}<e^{-3P}.
\]
Here $P=o(\sigma)$ and $\log(1/t)=O(\log\sigma)=o(P)$.
These exclusions give \eqref{eq:ambient-overlap-degree}.  They have not
been charged to the $o(n)$ exceptional set, as required.
\end{proof}

\begin{remark}\label{rem:geometry-validation}
The distinction between the two exceptional sets is essential.  The
second-moment validation argument can lose only $o(n)$ training points,
so it uses \eqref{eq:local-rich-lines}--\eqref{eq:overlap-degree}.
The stronger degree estimate \eqref{eq:ambient-overlap-degree} is used
only to bound the size of the predictor's ambient output, where an
additional $n\exp(CP)$ points are permitted.
\end{remark}
